\documentclass[10pt,a4paper]{amsart}
\usepackage{amsmath,amssymb,amsthm,graphicx}
% For commutative diagrams we use Xy-pic
\usepackage[all]{xy}

% We use 2cell for 2-commutative diagrams.
\xyoption{2cell}
% Theorem environments.
%
\theoremstyle{plain}
\newtheorem{theorem}[subsection]{Theorem}
\newtheorem{proposition}[subsection]{Proposition}
\newtheorem{lemma}[subsection]{Lemma}

\theoremstyle{definition}
\newtheorem{definition}[subsection]{Definition}
\newtheorem{example}[subsection]{Example}
\newtheorem{exercise}[subsection]{Exercise}
\newtheorem{situation}[subsection]{Situation}

\theoremstyle{remark}
\newtheorem{remark}[subsection]{Remark}
\newtheorem{remarks}[subsection]{Remarks}

\numberwithin{equation}{subsection}

% Macros
%
\def\lim{\mathop{\mathrm{lim}}\nolimits}
\def\colim{\mathop{\mathrm{colim}}\nolimits}
\def\Spec{\mathop{\mathrm{Spec}}}
\def\Hom{\mathop{\mathrm{Hom}}\nolimits}
\def\Ext{\mathop{\mathrm{Ext}}\nolimits}
\def\SheafHom{\mathop{\mathcal{H}\!\mathit{om}}\nolimits}
\def\SheafExt{\mathop{\mathcal{E}\!\mathit{xt}}\nolimits}
\def\Sch{\mathit{Sch}}
\def\Mor{\mathop{\mathrm{Mor}}\nolimits}
\def\Ob{\mathop{\mathrm{Ob}}\nolimits}
\def\Sh{\mathop{\mathit{Sh}}\nolimits}
\def\NL{\mathop{N\!L}\nolimits}
\def\CH{\mathop{\mathrm{CH}}\nolimits}
\def\proetale{{pro\text{-}\acute{e}tale}}
\def\etale{{\acute{e}tale}}
\def\QCoh{\mathit{QCoh}}
\def\Ker{\mathop{\mathrm{Ker}}}
\def\Im{\mathop{\mathrm{Im}}}
\def\Coker{\mathop{\mathrm{Coker}}}
\def\Coim{\mathop{\mathrm{Coim}}}

% Boxtimes
%
\DeclareMathSymbol{\boxtimes}{\mathbin}{AMSa}{"02}

%
% Macros for moduli stacks/spaces
%
\def\QCohstack{\mathcal{QC}\!\mathit{oh}}
\def\Cohstack{\mathcal{C}\!\mathit{oh}}
\def\Spacesstack{\mathcal{S}\!\mathit{paces}}
\def\Quotfunctor{\mathrm{Quot}}
\def\Hilbfunctor{\mathrm{Hilb}}
\def\Curvesstack{\mathcal{C}\!\mathit{urves}}
\def\Polarizedstack{\mathcal{P}\!\mathit{olarized}}
\def\Complexesstack{\mathcal{C}\!\mathit{omplexes}}
% \Pic is the operator that assigns to X its picard group, usage \Pic(X)
% \Picardstack_{X/B} denotes the Picard stack of X over B
% \Picardfunctor_{X/B} denotes the Picard functor of X over B
\def\Pic{\mathop{\mathrm{Pic}}\nolimits}
\def\Picardstack{\mathcal{P}\!\mathit{ic}}
\def\Picardfunctor{\mathrm{Pic}}
\def\Deformationcategory{\mathcal{D}\!\mathit{ef}}
\usepackage{cleveref}
\crefname{originalreference}{original reference}{original references}
\Crefname{originalreference}{Original reference}{Original references}
\makeatletter
\newlabel{lemma-describe-pullback}{{09XN}{source}{}{sourcepointer-1}{}}
\newlabel{lemma-describe-pullback@cref}{{[originalreference][0][]09XN}{[1][0][]source}{}{}{}}
\makeatother
\crefname{originalreference}{original reference}{original references}
\Crefname{originalreference}{Original reference}{Original references}
\makeatletter
\newlabel{lemma-proper-pushforward-stalk}{{0A3T}{source}{}{sourcepointer-2}{}}
\newlabel{lemma-proper-pushforward-stalk@cref}{{[originalreference][0][]0A3T}{[1][0][]source}{}{}{}}
\makeatother
\crefname{originallemma}{original lemma}{original lemmas}
\Crefname{originallemma}{Original lemma}{Original lemmas}
\makeatletter
\newlabel{topologies-lemma-characterize-sheaf}{{0DBQ}{source}{}{sourcepointer-3}{}}
\newlabel{topologies-lemma-characterize-sheaf@cref}{{[originallemma][0][]0DBQ}{[1][0][]source}{}{}{}}
\makeatother
\title{The pullback of a small etale sheaf to the ph-topology is described by the sections of its small pullback}
\author{The Stacks Project Authors}
\date{stacks-project:etale-cohomology.tex:local-snapshot-20261011; GNU FDL 1.2 or later}
\begin{document}
\maketitle
\noindent\textit{Editorial note.} All delivered mathematical text is the author's own, copied from the pinned Stacks Project chapter file etale-cohomology.tex; no statement or proof step was rewritten, reordered or removed.\par\medskip
\noindent\textit{Editorial note.} The counted proof is the author's complete argument: the proof is written out in full: the reduction to the ph-sheaf condition, the rewriting of the equalizer diagram through the two pushforwards, and the stalk computation showing that a section in the intersection of the two pullback images is constant on the whole product, hence comes from the stalk.\par\medskip
\noindent\textit{Editorial note.} The label lemma-describe-pullback-pi-ph names the counted statement, and the proof environment immediately following it in the source is delivered verbatim as its proof.\par\medskip
\noindent\textit{Editorial note.} The counted statement is carried by the author's own lemma environment in the chapter source.\par\medskip
\noindent\textit{Editorial note.} Pre-open gate: python operations/agent-authoring/tools/precard\_gate.py --source etale-cohomology.tex --label lemma-describe-pullback-pi-ph --start 22174 --end 22237 --local-id 3550 returned RESULT=OK (L3550 free, occupancy FREE); the delivered span is the pool's own [22174,22186] / [22188,22249], re-checked against the source boundaries.\par\medskip
\noindent\textit{Editorial note.} The statement and proof are copied from the pinned chapter file; the cited inputs are kept as external pointers with their canonical Stacks tags.\par\medskip
\noindent\textit{Allowed standard prerequisites (editorial).}\begin{itemize}
\item The description of the pullback in the etale topology. Same chapter, the lemma identifying pi\_S\textasciicircum{}\{-1\}F with the rule (f : X -> S) |-> Gamma(X, f\_small\textasciicircum{}\{-1\}F) for the etale topology. Supplies the etale part of the statement and the reduction of the ph-sheaf condition.
\item The characterisation of sheaves for the ph-topology. Topologies, the lemma reducing the ph-sheaf condition to an equalizer diagram for every proper surjective morphism. Supplies the equalizer criterion used in the proof.
\item Stalks of a proper pushforward as sections over the geometric fibre. Same chapter, the lemma describing the stalk of the pushforward along a proper morphism as the sections over the geometric fibre. Supplies the stalk computation at the end of the proof.
\end{itemize}\medskip
\ifdefined\hypertarget\hypertarget{sourcepointer-1}{}\fi
\noindent\textit{Source reference.} Original reference 09XN: sources/stacks-etale-cohomology/pinned-original-source.tex, source lines 4960--4960. Prerequisite pointer invoked at proof line(s) 22190; it names a standard fact of the same chapter that the counted argument uses, and every step of the counted proof built on it is retained verbatim in the delivered body.\par\medskip
\ifdefined\hypertarget\hypertarget{sourcepointer-2}{}\fi
\noindent\textit{Source reference.} Original reference 0A3T: sources/stacks-etale-cohomology/pinned-original-source.tex, source lines 18254--18254. Prerequisite pointer invoked at proof line(s) 22231; it names a standard fact of the same chapter that the counted argument uses, and every step of the counted proof built on it is retained verbatim in the delivered body.\par\medskip
\ifdefined\hypertarget\hypertarget{sourcepointer-3}{}\fi
\noindent\textit{Source reference.} Original lemma 0DBQ: sources/stacks-chow/tags.txt, source lines 16705--16705. Prerequisite pointer invoked at proof line(s) 22192; it names a standard fact of a sibling chapter that the counted argument uses, and every step of the counted proof built on it is retained verbatim in the delivered body.\par\medskip
\begin{lemma}
\label{lemma-describe-pullback-pi-ph}
With notation as above.
Let $\mathcal{F}$ be a sheaf on $S_\etale$. The rule
$$
(\Sch/S)_{ph} \longrightarrow \textit{Sets},\quad
(f : X \to S) \longmapsto \Gamma(X, f_{small}^{-1}\mathcal{F})
$$
is a sheaf and a fortiori a sheaf on $(\Sch/S)_\etale$.
In fact this sheaf is equal to
$\pi_S^{-1}\mathcal{F}$ on $(\Sch/S)_\etale$ and
$\epsilon_S^{-1}\pi_S^{-1}\mathcal{F}$ on $(\Sch/S)_{ph}$.
\end{lemma}

\begin{proof}
The statement about the \'etale topology is the content
of Lemma \ref{lemma-describe-pullback}. To finish the proof it
suffices to show that $\pi_S^{-1}\mathcal{F}$ is a sheaf for the ph
topology. By Topologies, Lemma \ref{topologies-lemma-characterize-sheaf}
it suffices to show that given a proper surjective morphism
$V \to U$ of schemes over $S$ we have an equalizer diagram
$$
\xymatrix{
(\pi_S^{-1}\mathcal{F})(U) \ar[r] &
(\pi_S^{-1}\mathcal{F})(V) \ar@<1ex>[r] \ar@<-1ex>[r] &
(\pi_S^{-1}\mathcal{F})(V \times_U V)
}
$$
Set $\mathcal{G} = \pi_S^{-1}\mathcal{F}|_{U_\etale}$.
Consider the commutative diagram
$$
\xymatrix{
V \times_U V \ar[r] \ar[rd]_g \ar[d] & V \ar[d]^f \\
V \ar[r]^f & U
}
$$
We have
$$
(\pi_S^{-1}\mathcal{F})(V) = \Gamma(V, f^{-1}\mathcal{G}) =
\Gamma(U, f_*f^{-1}\mathcal{G})
$$
where we use $f_*$ and $f^{-1}$ to denote functorialities between
small \'etale sites. Second, we have
$$
(\pi_S^{-1}\mathcal{F})(V \times_U V) =
\Gamma(V \times_U V, g^{-1}\mathcal{G}) =
\Gamma(U, g_*g^{-1}\mathcal{G})
$$
The two maps in the equalizer diagram come from the two maps
$$
f_*f^{-1}\mathcal{G} \longrightarrow g_*g^{-1}\mathcal{G}
$$
Thus it suffices to prove $\mathcal{G}$ is
the equalizer of these two maps of sheaves.
Let $\overline{u}$ be a geometric point of $U$. Set
$\Omega = \mathcal{G}_{\overline{u}}$.
Taking stalks at $\overline{u}$ by
Lemma \ref{lemma-proper-pushforward-stalk}
we obtain the two maps
$$
H^0(V_{\overline{u}}, \underline{\Omega}) \longrightarrow
H^0((V \times_U V)_{\overline{u}}, \underline{\Omega}) =
H^0(V_{\overline{u}} \times_{\overline{u}} V_{\overline{u}},
\underline{\Omega})
$$
where $\underline{\Omega}$ indicates the constant sheaf with value
$\Omega$. Of course these maps are the pullback by the projection maps.
Then it is clear that the sections coming from pullback
by projection onto the first factor are constant on the fibres of
the first projection, and sections coming from pullback
by projection onto the first factor are constant on the fibres of
the first projection. The sections in the intersection of the images
of these pullback maps are constant on all of
$V_{\overline{u}} \times_{\overline{u}} V_{\overline{u}}$, i.e.,
these come from elements of $\Omega$ as desired.
\end{proof}

\end{document}
